\documentclass[12pt]{article}
\usepackage{amsmath, amssymb, amsthm}
\usepackage{parskip}
\usepackage{geometry}
\usepackage{enumitem}
\usepackage{fancyhdr}

% Set margins to match a standard homework document
\geometry{a4paper, margin=1in}

\title{Recitation 6}
\parskip=1ex
\author{Math 311 Fall 2026 \\ Rutgers University}
\date{}

\begin{document}

\maketitle

\textbf{Exercise 1}
Let \( S \) be the set of positive integers whose
decimal representations contain no digit 0. Define
\[
    s_n = \sum_{1 \leq j \leq n} \frac{1}{j}
\]
for \( j \in \mathbb{S} \).
\begin{enumerate}[label=(\alph*)]
    \item For \( k \in \mathbb{N} \), let
          \[
              S_k = \{j \in \mathbb{S}: 10^{k-1} \leq  j \leq 10^k\}.
          \]
          Then every digit of \( j \) has 9 choices.
          Since there are \( k \) digits for \( S_{k} \), it follows that
          \[
              \vert S_k \vert = 9^k.
          \]
          By definition of \( S_k \),
          \[
              j \geq 10^{k-1},
          \]
          so
          \[
              \frac{1}{j} \leq \frac{1}{10^{k-1}}.
          \]
          Hence
          \[
              \begin{split}
                  \sum_{j \in S_k} \frac{1}{j} & \leq \frac{\vert S_k \vert }{10^{k-1}} \\
                                               & = 9 \left(\frac{9}{10}\right)^{k-1}.
              \end{split}
          \]

    \item Let \( \epsilon > 0 \). Since
          \[
              90 (\frac{9}{10})^{k-1} \to 0,
          \]
          we can choose \( K \in \mathbb{N} \) such that
          \[
              90 (\frac{9}{10})^{K-1} < \epsilon.
          \]
          Set \( N = 10^{K-1} \), and let \( m > n \geq N \).
          Then
          \[
              \begin{split}
                  \vert s_m - s_n \vert & = \sum_{n < j \leq m} \frac{1}{j}                          \\
                                        & \leq \sum_{k=K}^{\infty} \sum_{j \in S_k}                  \\
                                        & \leq \sum_{k=K}^{\infty} 9 \left(\frac{9}{10}\right)^{k-1} \\
                                        & = 90 \left(\frac{9}{10}\right)^{K-1}                       \\
                                        & < \epsilon.
              \end{split}
          \]
          So \( s_n \) is cauchy.
          Thus, the Cauchy Criterion says that \( s_n \to L \) for some \( L \in \mathbb{R} \).
          Since \( s_1 = 1 \) and \( S \) only contains positive terms, it follows that \( s_n \geq 1\) for all \( n \in \mathbb{N} \).
          Thus by the Order Limit Theorem, \( L \geq 1 > 0 \).
          Next, from part (a) we have
          \[
              \begin{aligned}
                  s_n
                   & \le \sum_{k=1}^{\infty}\sum_{j\in S_k}\frac1j         \\
                   & \le \sum_{k=1}^{\infty}9\left(\frac9{10}\right)^{k-1} \\
                   & =\frac9{1-\frac9{10}}                                 \\
                   & =90.
              \end{aligned}
          \]
          Therefore, by the Order Limit Theorem
          \[
              0 < L < 90.
          \]

\end{enumerate}

\textbf{Exercise 2}
\begin{enumerate}[label=(\alph*)]
    \item Let \( \{x_n\}_{n \in \mathbb{N}} \) be a real sequence, and define
          \[
              u_n = x_n + x_{n+1}, \qquad v_n = x_n + x_{n+2}.
          \]
          Suppose \( u_n \) and \( v_n \) are Cauchy. Notice that
          \[
              \begin{split}
                  u_n + u_{n-1} - v_{n-1} & = (x_n + x_{n+1}) + (x_{n-1} + x_n) - (x_{n-1} + x_{n+1}) \\
                                          & = 2x_n.
              \end{split}
          \]
          Let \( \epsilon > 0 \). Since \( u_n \) is cauchy,
          \[
              m,n \geq N_1 \implies \vert u_m - u_n \vert < \frac{2\epsilon}{3}.
          \]
          Similarly, since \( v_n \) is cauchy,
          \[
              m,n \geq N_2 \implies \vert v_m - v_n \vert < \frac{2\epsilon}{3}.
          \]
          Choose \( N = \max{\{N_1, N_2\}} + 1 \). Then
          \[
              \begin{split}
                  \vert x_m - x_{n} \vert & = \frac{1}{2} \vert (u_m + u_{m-1} - v_{m-1}) - (u_n + u_{n-1} - v_{n-1}) \vert                                      \\
                                          & \leq \frac{1}{2} \left(\vert u_m - u_n \vert + \vert u_{m-1} - u_{n-1} \vert + \vert v_{m-1} - v_{n-1} \vert \right) \\
                                          & \leq \frac{1}{2} \left(\frac{2\epsilon}{3} + \frac{2\epsilon}{3} + \frac{2\epsilon}{3}\right)                        \\
                                          & = \epsilon.
              \end{split}
          \]
          Therefore \( x_n \) is cauchy.

    \item Suppose \( \sum_{n=1}^{\infty} a_n^2 \) and \( \sum_{n=1}^{\infty} b_n^2 \) both converge.
          Notice that
          \[
              \begin{split}
                  (\vert a_n \vert - \vert b_n \vert )^2            & \geq  0                                \\
                  a_n^2 - 2 \vert a_n \vert \vert b_n \vert + b_n^2 & \geq 0                                 \\
                  \frac{a_n^2 + b_n^2}{2}                           & \geq \vert a_n \vert \vert b_n \vert .
              \end{split}
          \]
          Next, by the Algebraic Limit Theorem, we have that
          \[
              \sum_{n=1}^{\infty} \frac{a_n^2 + b_n^2}{2} \quad \text{converges.}
          \]
          Moreover, by the comparison test
          \[
              \sum_{n=1}^{\infty} \vert a_n \vert \vert b_n \vert \quad \text{converges.}
          \]
          Hence \( \sum_{n=1}^{\infty}  a_n b_n \) converges absolutely.
          Then, by theorem 2.7.10, we can conclude that ever rearrangement of this series also converges.
\end{enumerate}

\textbf{Exercise 3}
For \( m,n \in \mathbb{N} \), define
\[
    a_{m,n} = \frac{mn^2}{m^2 + n^4}.
\]
\begin{enumerate}[label=(\alph*)]
    \item For \( \lim_{m\to\infty}\lim_{n\to\infty}a_{m,n}\), we first divide through by \( n^4 \) and then take \( n \to \infty \),
          \[
              \lim_{n \to \infty} a_{m,n} = \lim_{n \to \infty} \frac{\frac{m}{n^2}}{\frac{m^2}{n^4} + 1} = 0.
          \]
          Taking \( m \to \infty \) does not do anything, so \( \lim_{m\to\infty}\lim_{n\to\infty}a_{m,n} = 0 \).

          For \( \lim_{n\to\infty}\lim_{m\to\infty}a_{m,n}\), we first divide through by \( m^2 \) and then take \( m \to \infty \),
          \[
              \lim_{m \to \infty} a_{m,n} = \lim_{n \to \infty} \frac{\frac{n^2}{m}}{1 + \frac{n^4}{m^2}} = 0.
          \]
          Taking \( n \to \infty \) does not do anything, so \( \lim_{n\to\infty}\lim_{m\to\infty}a_{m,n} = 0 \).

    \item No, consider \( m = n^2 \).
          Then
          \[
              \begin{aligned}
                  a_{n^2,n}
                   & =\frac{n^2\cdot n^2}{(n^2)^2+n^4} \\
                   & =\frac{n^4}{2n^4}                 \\
                   & =\frac12.
              \end{aligned}
          \]

          Hence,
          \[
              \lim_{n\to\infty}a_{n^2,n}=\frac12\neq0.
          \]

    \item Suppose that \(a_{m,n}\to L\) as \(m,n\to\infty\).
          Because \(a_{n,n}\to0\), the uniqueness of limits implies that \(L=0\). Moreover, convergence of the double sequence to \(L\) forces its diagonal sequence to converge to \(L\).
          Now choose
          \[
              \boxed{\varepsilon_0=\frac14}.
          \]

          Let \(N\in\mathbb N\) be arbitrary. Choose \(n\ge N\) and set \(m=n^2\). Since \(n\ge1\),
          \[
              m=n^2\ge n\ge N.
          \]

          But
          \[
              \begin{aligned}
                  |a_{m,n}-L|
                   & =|a_{n^2,n}-0|            \\
                   & =\left|\frac12\right|     \\
                   & =\frac12                  \\
                   & \ge\frac14=\varepsilon_0.
              \end{aligned}
          \]

          Thus, for every \(N\in\mathbb N\), there exist \(m,n\ge N\) satisfying
          \[
              |a_{m,n}-L|\ge\varepsilon_0.
          \]

          This contradicts the definition of double-sequence convergence.
          Therefore,
          \[
              \{a_{m,n}\}_{m,n\in\mathbb N}\text{ does not converge as }m,n\to\infty.
          \]
\end{enumerate}




\end{document}

